\frametitle{The Error Splits, and What the Split Assigns}
\footnotesize

Every algebraic error splits uniquely and $A$-orthogonally,
\[
e=e_C+e_F,\qquad e_C=\PiC e\in\operatorname{range}(P),\qquad e_F=\PiF e ,
\]
characterised by three identities, all of which hold exactly:
\[
\boxed{\;e_C^{\mathsf T}A\,e_F=0,\qquad P^{\mathsf T}A\,e_F=0,\qquad
\normA{e}^2=\normA{e_C}^2+\normA{e_F}^2\;}
\]
The \emph{second} identity carries the interpretation. Since
$\PiF=I-P(P^{\mathsf T}AP)^{-1}P^{\mathsf T}A$, applying $P^{\mathsf T}A$ to $e_F$ gives
zero: $e_F$ is exactly the part of the error the coarse grid \alert{cannot see}. The
correction, which acts by $PA_c^{-1}P^{\mathsf T}$ on the residual, is identically blind to
it --- equivalently, the residual produced by $e_F$ is orthogonal to the coarse space.

\begin{columns}[T]
\begin{column}{0.47\textwidth}
$\R^n=\operatorname{range}(AP)\oplus\ker(P^{\mathsf T})$, of dimensions $n_c$ and
$n-n_c$. A complement residual satisfies $P^{\mathsf T}r_F=0$; a coarse residual lies in
$\operatorname{range}(AP)$; and the two subspaces meet only at the origin. So a step trained
to act on $e_F$ \emph{has never been shown a single input from
$\operatorname{range}(AP)$}, and its behaviour there is a pure extrapolation --- whether
that is harmless or harmful is a measurement, not an assumption.
\end{column}
\begin{column}{0.47\textwidth}
So the smoother's duty is not ``reduce the error''. It is narrower: $e_C$ is the
coarse-grid correction's business, and the correction is identically blind to $e_F$, so
\alert{if the smoother does not reduce $e_F$ then nothing in the cycle does}. That is what
makes the first experiment an \emph{extrapolation test} rather than a memorisation test,
and it is the vocabulary in which the defect found there is stated. The classical
smoothers are the right comparison set: they too are judged by what they do to the whole
space, and the split says which part of that matters.
\end{column}
\end{columns}
